% TOP_PROBLEM: {"id": 2621, "title": "Kourovka Notebook Problem 21.112", "category_id": 17, "category": {"id": 17, "name": "group_theory", "display_name": "Group Theory", "description": "Problems about groups, group actions, representations, and related algebraic structures.", "slug": "group-theory", "order_index": 17, "created_at": "2026-07-31T15:26:25.671Z"}, "status": "open", "problem_number": "KOU-21.112", "set_id": 10}
% FIRST_POSTED: 2026-10-01
% ENTRY_KIND: statement_only
% UnsolvedMath ID 2621; source catalogue code KOU-21.112
% !TeX program = lualatex
% Definition and sources only; this file is not an LLM solution attempt.
\documentclass[11pt,a4paper]{article}
\usepackage[margin=25mm]{geometry}
\usepackage{fontspec}
\setmainfont{lmroman10-regular.otf}[BoldFont=lmroman10-bold.otf,
 ItalicFont=lmroman10-italic.otf,BoldItalicFont=lmroman10-bolditalic.otf]
\usepackage{amsmath,amssymb,mathtools}
\usepackage{unicode-math}
\setmathfont{latinmodern-math.otf}
\AtBeginDocument{\let\setminus\smallsetminus}
\usepackage{xurl,hyperref}
\hypersetup{colorlinks=true,urlcolor=blue}
\setcounter{secnumdepth}{0}
\setlength{\parindent}{0pt}
\setlength{\parskip}{5pt}
\setlength{\emergencystretch}{3em}
\begin{document}
\section*{Kourovka Notebook Problem 21.112}\label{2621}
{\small UnsolvedMath ID 2621 · KOU-21.112 · Group Theory · Kourovka Notebook - New Problems, Issue 21}
\subsection{Definitions and mathematical statement}
A complete class $\mathcal X$ is a nonempty class of finite groups closed under
subgroups, quotient groups, and extensions: if $N\trianglelefteq G$ and
$N,G/N\in\mathcal X$, then $G\in\mathcal X$.
Assume $\mathcal X$ is proper, so it does not contain every finite group.
Its symmetric boundary is $n=\max\{r:S_r\in\mathcal X\}$, where $S_r$ is
the permutation group on $r$ points. This maximum is finite because every
finite group embeds in some symmetric group.

The Baer--Suzuki width $\operatorname{BS}(\mathcal X)$ is the least
nonnegative integer $m$ with the following property: for every finite group
$G$ and each conjugacy class $D\subset G$, if
$\langle d_1,\ldots,d_m\rangle\in\mathcal X$ for every tuple
$(d_1,\ldots,d_m)\in D^m$, then $\langle D\rangle\in\mathcal X$.
Repetitions in the tuple are allowed. The source establishes finiteness of
this width and bounds it in terms of the symmetric boundary.

For each positive integer $n\ne3$, set
\[
f_+(n)=\max_{\partial_s\mathcal X=n}\operatorname{BS}(\mathcal X),\qquad
f_-(n)=\min_{\partial_s\mathcal X=n}\operatorname{BS}(\mathcal X),
\]
where $\partial_s$ denotes symmetric boundary and the extrema range over all
proper complete classes. Such classes exist for precisely the indicated $n$.
Determine $f_+(4),f_+(5),f_+(6)$, and decide whether $f_-(n)=n$ for every
positive $n\ne3$.
\subsection{Short English statement}
Find the largest Baer--Suzuki widths at symmetric boundaries four, five and six, and decide whether the smallest width always equals the boundary.
\subsection{Sources}
\begin{itemize}
\item[S1] ULAM.AI and the UnsolvedMath Contributors, supplied problems.json, record 2621 (KOU-21.112), Kourovka Notebook - New Problems, Issue 21. Catalogue identity and source transcription; CC BY 4.0.
\url{https://huggingface.co/datasets/ulamai/UnsolvedMath}.
\item[S2] The Kourovka Notebook, 21st edition, new problems, Problem 21.112. Expanded formulation of the supplied catalogue statement.
\url{https://arxiv.org/pdf/1401.0300v47}.
\end{itemize}
\end{document}
